\documentclass[11pt,a4paper]{article}
\usepackage[margin=23mm,headheight=14pt,headsep=8.5mm,footskip=10mm]{geometry}
\usepackage{amsmath,amssymb}
\usepackage{assets/tex/cm-bea}
\usepackage{xeCJK}
\setmainfont{Latin Modern Roman}
\setCJKmainfont{FandolSong-Regular.otf}[Path=assets/fonts/fandol/,BoldFont=FandolSong-Bold.otf,ItalicFont=FandolSong-Regular.otf]
\setCJKsansfont{FandolSong-Regular.otf}[Path=assets/fonts/fandol/,BoldFont=FandolSong-Bold.otf]
\xeCJKsetup{PunctStyle=quanjiao,CheckSingle=true}
\usepackage{setspace}\setstretch{1.12}
\setlength{\parindent}{2em}\setlength{\parskip}{3pt}
\usepackage{titlesec}
\titleformat{\section}{\centering\normalfont\large\bfseries}{}{0pt}{}
\titlespacing*{\section}{0pt}{1.5\baselineskip}{0.8\baselineskip}
\usepackage{needspace,enumitem,longtable,booktabs,array,calc}
\providecommand{\tightlist}{\setlength{\itemsep}{3pt}\setlength{\parskip}{0pt}}
\usepackage{fancyhdr}\pagestyle{fancy}\fancyhf{}
\fancyhead[L]{\small 数学手记}\fancyhead[R]{\small 楼红卫 B 组 · 第四十二批}
\fancyfoot[C]{\thepage}\renewcommand{\headrulewidth}{0.4pt}
\usepackage[hidelinks,unicode,linktoc=all]{hyperref}\usepackage{bookmark}
\hypersetup{pdftitle={楼红卫数学分析 B289—B295：曲线曲面积分、行列式恒等式与循环常数},pdfauthor={Vincent Arqain}}
\renewcommand{\contentsname}{题号目录}\setcounter{tocdepth}{1}\setcounter{secnumdepth}{0}
\widowpenalty=10000\clubpenalty=10000\displaywidowpenalty=10000
\predisplaypenalty=10000\interdisplaylinepenalty=10000
\raggedbottom\emergencystretch=2em
\begin{document}
\begin{center}
{\Large\bfseries 楼红卫《数学分析》B 组习题解答\par}
\vspace{7pt}{\large 第四十二批\quad B289—B295\par}
\vspace{5pt}{曲线曲面积分、行列式恒等式与循环常数\par}
\end{center}
本批收录习题 11.1—11.4 各节 B 组第 1 题及 11.5.B 第 1—3 题。正文依次给出原题、基础知识、思路来源、完整证明、易错条件与方法总结。数学公式采用 CM bea，中文采用 Fandol 宋体；公式编号与网页一致。
\begingroup\setlength{\parskip}{0pt}\tableofcontents\endgroup

\Needspace{9\baselineskip}
\section{问题 B289}\label{ux95eeux9898-b289}

\textbf{原题号：} 习题 11.1.B 第 1
题。\href{https://zj139.github.io/math-notes/files/problems/楼分析习题-v2.pdf\#page=34}{原题，PDF
第 34 页（印刷第 32 页）}

\textbf{题目}　试将定积分的一些性质移植到第一型曲线积分。

\Needspace{3\baselineskip}\textbf{基础知识}　分段正则曲线的弧长参数；定积分的线性、可加性、保序性、中值定理及一致收敛下的逐项积分。

\Needspace{3\baselineskip}\textbf{思路来源}　第一型曲线积分以弧长为正的积分权。用弧长参数把曲线拉回一个区间，定积分性质便有了共同的证明入口；但方向、连通性与重复走过曲线的次数必须说明。

\Needspace{3\baselineskip}\textbf{解}　设 \(C\) 为有限长、连通、分段 \(C^1\) 正则曲线，长度为
\(L>0\)，先约定只走过一次。取弧长参数表示 \(r:[0,L]\to C\)，则对可积的
\(f\)，

\begin{equation}
\int_C f\,ds=\int_0^L f(r(s))\,ds.
\tag{1}
\end{equation}

以下各积分均假设存在；论及实数大小时函数取实值。

\textbf{线性与可加性。} 对常数 \(\alpha,\beta\)，由（1）有

\begin{equation}
\int_C(\alpha f+\beta g)\,ds
=\alpha\int_Cf\,ds+\beta\int_Cg\,ds.
\tag{2}
\end{equation}

若把参数区间在有限个点处分段，相应的曲线段为 \(C_1,\ldots,C_k\)，则
\(\int_C f\,ds=\sum_{j=1}^k\int_{C_j}f\,ds\)。端点的重复不影响积分；若几何上某一段被重复走过，则必须计入其遍历重数，不能只按点集取并。

\textbf{保序与绝对值估计。} 由正权 \(ds\)，\(f\le g\) 蕴涵
\(\int_C f\,ds\le\int_Cg\,ds\)，特别地

\begin{equation}
\begin{aligned}
\int_C1\,ds&=L,\\
m\le f\le M&\ \Longrightarrow\ mL\le\int_Cf\,ds\le ML,\\
\left|\int_Cf\,ds\right|&\le\int_C|f|\,ds
\le L\sup_C|f|.
\end{aligned}
\tag{3}
\end{equation}

最后一行对复值函数也成立：若 \(I=\int_Cf\,ds\ne0\)，乘以
\(\overline I/|I|\) 后取实部，即得 \(|I|\le\int_C|f|\,ds\)。若 \(f\ge0\)
且在 \(C\) 上连续，则积分为零还推出 \(f\equiv0\)；否则在某个弧长小区间上
\(f\) 有正下界，积分必为正。没有连续性时，只能推出 \(f=0\)
几乎处处（关于弧长）。

\textbf{积分中值定理。} 设 \(f\) 在 \(C\) 上连续，\(g\ge0\) 可积且
\(\int_Cg\,ds>0\)。记 \(m=\min_C f\)、\(M=\max_Cf\)，则

\begin{equation}
m\le\frac{\int_Cfg\,ds}{\int_Cg\,ds}\le M.
\tag{4}
\end{equation}

由于 \(C\) 连通，连续函数 \(f\) 取遍 \([m,M]\)，故存在 \(\xi\in C\)，使

\begin{equation}
\int_Cfg\,ds=f(\xi)\int_Cg\,ds.
\tag{5}
\end{equation}

取 \(g=1\) 即得
\(\int_Cf\,ds=f(\xi)L\)。连通性不能去掉：两条互不相交、等长的线段上分别取
\(f=0\)、\(f=1\)，平均值为 \(1/2\)，却没有任何点取到它。

\textbf{Hölder 与 Cauchy--Schwarz 不等式。} 对
\(p,q>1\)、\(1/p+1/q=1\)，把区间上的 Hölder 不等式用于 \(f\circ r\) 与
\(g\circ r\)，得到

\begin{equation}
\int_C|fg|\,ds\le
\left(\int_C|f|^p\,ds\right)^{1/p}
\left(\int_C|g|^q\,ds\right)^{1/q}.
\tag{6}
\end{equation}

取 \(p=q=2\) 即为 Cauchy--Schwarz
不等式。这里要求右端两个幂积分有限，才能得到有限值估计。

\textbf{一致收敛与积分交换。} 若 \(f_n\) 在 \(C\) 上可积且一致收敛于
\(f\)，则由区间上的相应定理 \(f\) 可积，并且

\begin{equation}
\left|\int_Cf_n\,ds-\int_Cf\,ds\right|
\le L\sup_C|f_n-f|\longrightarrow0.
\tag{7}
\end{equation}

同理，一致收敛的函数项级数可以逐项作弧长积分。在 Lebesgue
积分的框架下，单调收敛和控制收敛定理也可由（1）直接移植；控制函数必须关于弧长可积。

\textbf{方向与参数无关性。} 若 \(r=r(t)\) 是任意分段正则参数表示，则
\(ds=|r'(t)|dt\)。反向参数用速度的模，故

\begin{equation}
\int_{-C}f\,ds=\int_Cf\,ds.
\tag{8}
\end{equation}

一般的单调正则重参数化也由一元换元公式得到同一个值。因此第一型曲线积分不依赖方向，不能把定积分``交换上下限变号''的规则照搬为（8）的负号。

\Needspace{3\baselineskip}\textbf{易错条件}　中值点的存在需要连通性及实值连续函数；积分为零不能对任意可积函数推出逐点为零；反向不变不等于重复遍历不变。无限长曲线不能直接使用式（7）中的有限常数
\(L\)。

\Needspace{3\baselineskip}\textbf{方法总结}　先把几何对象参数化为区间，再识别正权
\(ds\)。线性、保序与极限运算继承自区间积分；方向和遍历重数则由定义本身决定。

\Needspace{9\baselineskip}
\section{问题 B290}\label{ux95eeux9898-b290}

\textbf{原题号：} 习题 11.2.B 第 1
题。\href{https://zj139.github.io/math-notes/files/problems/楼分析习题-v2.pdf\#page=35}{原题，PDF
第 35 页（印刷第 33 页）}

\textbf{题目}　设 \(a,b,c\)
为实数，\(A=\sqrt{a^2+b^2+c^2}\)，\(f\in L^1[-A,A]\)，证明
Poisson（泊松）公式

\begin{equation}
\iint_{x^2+y^2+z^2=1}f(ax+by+cz)\,dS
=2\pi\int_{-1}^1f(Au)\,du.
\tag{9}
\end{equation}

\Needspace{3\baselineskip}\textbf{基础知识}　正交变换保持球面与面积；曲面参数化的面积因子；绝对可积函数的换元与
Fubini 定理。

\Needspace{3\baselineskip}\textbf{思路来源}　把向量 \((a,b,c)\)
旋转到坐标轴方向，线性函数便只依赖球面的高度。单位球面的等高薄带面积为
\(2\pi\,du\)，这是公式中常数 \(2\pi\) 的来源。

\Needspace{3\baselineskip}\textbf{证明}　先设 \(A>0\)。选取正交矩阵 \(O\)，使其第一行为
\((a,b,c)/A\)。令 \((u,v,w)^T=O(x,y,z)^T\)，则 \(ax+by+cz=Au\)，球面变为
\(u^2+v^2+w^2=1\)。正交变换保持切向量的 Gram 行列式，因而保持面积元。

在去掉两极和一条经线后的球面上，取参数表示

\begin{equation}
X(u,\theta)=\big(u,\sqrt{1-u^2}\cos\theta,
\sqrt{1-u^2}\sin\theta\big),
\quad -1<u<1,\quad0<\theta<2\pi.
\tag{10}
\end{equation}

直接计算得到

\begin{equation}
|X_u|^2=\frac1{1-u^2},\qquad
|X_\theta|^2=1-u^2,\qquad
X_u\cdot X_\theta=0,\qquad
|X_u\times X_\theta|=1.
\tag{11}
\end{equation}

被去掉的经线和两极面积为零。先对绝对值使用面积公式，得到

\begin{equation}
\iint_{S^2}|f(ax+by+cz)|\,dS
=2\pi\int_{-1}^1|f(Au)|\,du
=\frac{2\pi}{A}\int_{-A}^A|f(t)|\,dt<\infty.
\tag{12}
\end{equation}

故换序合法，再对 \(f\) 本身计算，

\begin{equation}
\iint_{S^2}f(ax+by+cz)\,dS
=\int_{-1}^1\int_0^{2\pi}f(Au)\,d\theta\,du
=2\pi\int_{-1}^1f(Au)\,du.
\tag{13}
\end{equation}

若把 \(L^1\) 看作几乎处处相等的等价类，式（12）还说明：当 \(A>0\)
时，改变 \(f\) 在一维零测集上的值不会改变球面积分。

若 \(A=0\)，则 \(a=b=c=0\)。在给定有限值 \(f(0)\)
的通常逐点解释下，两边均为 \(4\pi f(0)\)，公式仍成立。但此时不能除以
\(A\)；若题目仅给出退化区间 \([0,0]\) 上的 \(L^1\)
等价类，该等价类不确定
\(f(0)\)，必须另行指定这个点值才能解释两边。它是定义上的退化情形，不是公式在
\(A>0\) 时的附加限制。

\Needspace{3\baselineskip}\textbf{易错条件}　不能只对连续函数作形式换元就忽略 \(L^1\)
假设；式（12）保证复合函数确实可积。球面极点处参数退化不影响积分，但应说明被遗漏集合面积为零。

\Needspace{3\baselineskip}\textbf{方法总结}　旋转对称性消去线性函数的方向，等高参数消去面积因子。先对绝对值建立可积性，再处理有符号或复值函数，能同时解决换序和代表元问题。

\Needspace{9\baselineskip}
\section{问题 B291}\label{ux95eeux9898-b291}

\textbf{原题号：} 习题 11.3.B 第 1
题。\href{https://zj139.github.io/math-notes/files/problems/楼分析习题-v2.pdf\#page=35}{原题，PDF
第 35---36 页（印刷第 33---34 页）}

\textbf{题目}　对于复平面中的 \(C^1\) 曲线 \(C\) 及其上的复值连续函数
\(f\)，记

\begin{equation}
f(z)\,dz=f(x+iy)(dx+i\,dy)
=f(x+iy)\,dx+i f(x+iy)\,dy.
\tag{14}
\end{equation}

若 \(s\) 为弧长参数，\(C_{s,\Delta s}\) 为弧长从 \(s\) 到 \(s+\Delta s\)
的那一段，则

\begin{equation}
\lim_{\Delta s\to0}
\frac{\left|\int_{C_{s,\Delta s}}dz\right|}{|\Delta s|}=1.
\tag{15}
\end{equation}

这相当于 \(|dz|=ds\)。证明

\begin{equation}
\left|\int_C f(z)\,dz\right|\le\int_C|f(z)|\,ds.
\tag{16}
\end{equation}

\Needspace{3\baselineskip}\textbf{基础知识}　复线积分的参数定义；正则曲线的单位切向量；复积分的三角不等式。

\Needspace{3\baselineskip}\textbf{思路来源}　\(dz\) 含有切向方向，而 \(ds\)
只记录长度。把曲线积分写成 \(f(z(t))z'(t)\)
的一元积分，模长估计便把方向信息去掉；不能把``积分的模''与``模的积分''视作相等。

\Needspace{3\baselineskip}\textbf{证明}　取参数表示 \(z=z(t)\)，\(t\in[\alpha,\beta]\)。按定义，

\begin{equation}
\int_C f(z)\,dz=\int_\alpha^\beta f(z(t))z'(t)\,dt,
\qquad
\int_C|f(z)|\,ds=\int_\alpha^\beta|f(z(t))|\,|z'(t)|\,dt.
\tag{17}
\end{equation}

令 \(I\) 为左端的复积分。若 \(I=0\)，结论显然；否则取
\(\lambda=\overline I/|I|\)，则 \(|\lambda|=1\)，由实部的线性有

\begin{equation}
\begin{aligned}
|I|&=\operatorname{Re}(\lambda I)
=\int_\alpha^\beta\operatorname{Re}\big(\lambda f(z(t))z'(t)\big)\,dt\\
&\le\int_\alpha^\beta|f(z(t))z'(t)|\,dt
=\int_C|f(z)|\,ds.
\end{aligned}
\tag{18}
\end{equation}

这证明（16）；参数导数偶有零点时，（17）---（18）仍然有效。

再说明题中的（15）。在正则的弧长参数处，\(|z'(s)|=1\)，而

\begin{equation}
\int_{C_{s,\Delta s}}dz=z(s+\Delta s)-z(s),\qquad
\lim_{\Delta s\to0}\frac{|z(s+\Delta s)-z(s)|}{|\Delta s|}
=|z'(s)|=1.
\tag{19}
\end{equation}

所以（15）是``短弦长与短弧长之比趋于一''，并非有限弧段上弦长等于弧长。对分段正则曲线，逐段证明（16）后用三角不等式相加即可。

还可读出等号条件：若右端非零，设 \(T(s)=z'(s)\)
为单位切向量，则取等号当且仅当存在固定的 \(|\lambda|=1\)，使

\begin{equation}
\lambda f(z(s))T(s)\in[0,\infty)
\quad\text{对几乎所有 }s\text{ 成立}.
\tag{20}
\end{equation}

这是因为（18）中的差为非负函数，其积分为零恰好要求它几乎处处为零；而
\(\operatorname{Re}w=|w|\) 当且仅当 \(w\) 为非负实数。

\Needspace{3\baselineskip}\textbf{易错条件}　``\(C^1\)
参数曲线''若未要求正则，不能不加说明地声称每点都有光滑的单位速参数；不过不等式本身仍由（17）成立。\(|dz|=ds\)
不能读成 \(|\int dz|=\int ds\)。

\Needspace{3\baselineskip}\textbf{方法总结}　复积分估计可用一个固定相位把最终积分旋转到正实轴，再逐点取实部比较。这同时给出最优的模估计及等号成立条件。

\Needspace{9\baselineskip}
\section{问题 B292}\label{ux95eeux9898-b292}

\textbf{原题号：} 习题 11.4.B 第 1
题。\href{https://zj139.github.io/math-notes/files/problems/楼分析习题-v2.pdf\#page=36}{原题，PDF
第 36 页（印刷第 34 页）}

\textbf{题目}　若请你对 \(\mathbb R^3\) 中的曲面 \(\Sigma\) 及函数 \(P\)
定义 \(\iint_\Sigma P(x,y,z)\,dy\,dz\)，则对于 \(\Sigma\) 和 \(P\)
的要求可以怎么提？

\Needspace{3\baselineskip}\textbf{基础知识}　有向曲面及单位法向量；正则曲面片；参数变换的
Jacobian；第二型曲面积分与通量的关系。

\Needspace{3\baselineskip}\textbf{思路来源}　\(dy\,dz\)
是带方向的投影面积，不能只由无向的点集决定。先给一组容易验证的充分条件，再用参数变换证明定义一致，最后指出哪些条件可以放宽。

\Needspace{3\baselineskip}\textbf{解}　一组经典的充分条件是：\(\Sigma\) 为紧致、可定向、分片
\(C^1\)
正则曲面，已选定一侧；能分成有限个面积有限的正则曲面片，片间只在面积为零的边界上重叠；\(P\)
为 \(\Sigma\)
上的实值连续函数。曲面允许有边界，也不要求闭合。记所选单位法向量为
\(\nu=(\nu_1,\nu_2,\nu_3)\)，定义

\begin{equation}
\iint_\Sigma P\,dy\,dz:=\iint_\Sigma P\nu_1\,dS.
\tag{21}
\end{equation}

右端为第一型曲面积分。\(P\)
有界，\(|\nu_1|\le1\)，曲面面积有限，故此积分存在。

对一个与选定方向相容的正则参数片 \(X(u,v)=(x(u,v),y(u,v),z(u,v))\)，有

\begin{equation}
\begin{aligned}
\nu&=\frac{X_u\times X_v}{|X_u\times X_v|},\\
\nu_1\,dS&=(y_u z_v-y_v z_u)\,du\,dv,\\
\iint_{X(D)}P\,dy\,dz
&=\iint_D P(X(u,v))(y_u z_v-y_v z_u)\,du\,dv.
\end{aligned}
\tag{22}
\end{equation}

若另一个相容参数满足 \((u,v)=h(\xi,\eta)\)，链式法则给出

\begin{equation}
X_\xi\times X_\eta
=\det Dh\,(X_u\times X_v).
\tag{23}
\end{equation}

相容方向意味着
\(\det Dh>0\)，由二重积分换元公式，（22）的值不变；改用反向参数而不同时修正方向则变号。有限分片上的定义由可加性拼合，公共边界面积为零，不产生重复计数。这就证明（21）不是依赖某个偶然参数的定义。

这些条件是充分条件，并非最弱条件。若采用 Lebesgue 曲面积分，\(P\)
连续可放宽为可测且

\begin{equation}
\iint_\Sigma |P\nu_1|\,dS<\infty.
\tag{24}
\end{equation}

较易检验的 \(P\in L^1(\Sigma,dS)\)
也足够，但比（24）更强。非紧曲面可以用局部有限的参数片拼合，仍以（24）保证绝对收敛与分片次序无关。若允许自交参数曲面，须事先约定各片的方向与遍历重数，否则仅给出几何像集不足以确定积分。

不需要 \(P\) 有偏导数，不需要 \(\Sigma\) 是整个 \(yz\)
平面上的单值图像，也不需要 \(\nu_1\) 处处非零。例如 \(z=0\)
上的一块平面片，其
\(\nu_1=0\)，故此分量的积分恒为零。若曲面不可定向，则一般不能为任意
\(P\)
赋予一致的全局有向积分；应改为指定的有向曲面片之和，而不能忽略方向选择。

\Needspace{3\baselineskip}\textbf{易错条件}　投影可能重叠，不能把所有点简单投影到 \(yz\)
平面再只算一次。取绝对值 \(|\nu_1|\)
会得到无符号的投影面积权，已经改变了第二型积分；曲面光滑也不自动保证已选定全局方向。

\Needspace{3\baselineskip}\textbf{方法总结}　先用 \(P\nu_1dS\) 给出几何定义，再由参数 Jacobian
检查其一致性。把``定义积分所需的可积性''与``应用 Gauss
公式所需的可微性''分开，能避免提出多余条件。

\Needspace{9\baselineskip}
\section{问题 B293}\label{ux95eeux9898-b293}

\textbf{原题号：} 习题 11.5.B 第 1
题。\href{https://zj139.github.io/math-notes/files/problems/楼分析习题-v2.pdf\#page=37}{原题，PDF
第 37 页（印刷第 35 页）}

\textbf{题目}　设 \(a,b,c>0\)，\(\Sigma\) 为椭球面
\(x^2/a^2+y^2/b^2+z^2/c^2=1\)。利用 Gauss 公式计算

\begin{equation}
I=\iint_\Sigma
\frac{dS}{(x^2+y^2+z^2)^{3/2}
\sqrt{x^2/a^4+y^2/b^4+z^2/c^4}}.
\tag{25}
\end{equation}

\Needspace{3\baselineskip}\textbf{基础知识}　水平曲面的外单位法向量；Gauss
公式；挖去奇点后的边界方向。

\Needspace{3\baselineskip}\textbf{思路来源}　根号恰好是椭球法向量归一化的因子，而椭球方程使
\((x,y,z)\) 与未归一化法向量的内积等于一。因此被积式可识别为向量场
\((x,y,z)/(x^2+y^2+z^2)^{3/2}\) 的外向通量。

\Needspace{3\baselineskip}\textbf{解}　记
\(r=\sqrt{x^2+y^2+z^2}\)、\(H=\sqrt{x^2/a^4+y^2/b^4+z^2/c^4}\)。在椭球面上，外单位法向量为

\begin{equation}
\nu=\frac{(x/a^2,y/b^2,z/c^2)}{H}.
\tag{26}
\end{equation}

令 \(F=(x,y,z)/r^3\)（\(r>0\)），则由椭球方程，

\begin{equation}
F\cdot\nu=\frac{x^2/a^2+y^2/b^2+z^2/c^2}{r^3H}
=\frac1{r^3H},\qquad
\operatorname{div}F=\frac3{r^3}-\frac{3r^2}{r^5}=0.
\tag{27}
\end{equation}

设 \(E\) 为椭球内部。取 \(0<\varepsilon<\min(a,b,c)\)，则闭球
\(\overline B_\varepsilon(0)\) 包含在 \(E\) 内。在
\(E\setminus\overline B_\varepsilon(0)\) 上应用 Gauss 公式，\(F\)
在该区域闭包上光滑；内边界相对于挖孔区域的外单位法向量为
\(-(x,y,z)/\varepsilon\)，于是

\begin{equation}
\begin{aligned}
0&=\iint_\Sigma F\cdot\nu\,dS
-\iint_{\partial B_\varepsilon}
F\cdot\frac{(x,y,z)}{\varepsilon}\,dS,\\
I&=\iint_{\partial B_\varepsilon}\frac1{\varepsilon^2}\,dS
=\frac{4\pi\varepsilon^2}{\varepsilon^2}=4\pi.
\end{aligned}
\tag{28}
\end{equation}

每个允许的 \(\varepsilon\)
都给出同一结果，不需要再交换极限与积分。原积分为第一型曲面积分，本身无方向；外向法线只是在将它改写为通量时选用。

\Needspace{3\baselineskip}\textbf{易错条件}　\(F\)
在原点有奇性，不能直接在整个椭球内部套用``散度为零''，否则会错误地得到零。挖孔后的内边界方向必须反向，不能与小球本身的外向法线混淆。

\Needspace{3\baselineskip}\textbf{方法总结}　先从曲面方程识别法向量，再把复杂的第一型积分写成简单场的通量。对孤立奇点挖孔，将任意包围面上的通量转移到球面，是
Gauss 公式的典型构造。

\Needspace{9\baselineskip}
\section{问题 B294}\label{ux95eeux9898-b294}

\textbf{原题号：} 习题 11.5.B 第 2
题。\href{https://zj139.github.io/math-notes/files/problems/楼分析习题-v2.pdf\#page=37}{原题，PDF
第 37 页（印刷第 35 页）}

\textbf{题目}　设 \(A\) 是 \(n\) 阶方阵，\(\Omega\) 为 \(\mathbb R^n\)
中具有 \(C^1\)
边界的有界区域，\(\varphi\in C_0^2(\Omega;\mathbb R^n)\)。证明

\begin{equation}
\int_\Omega\det(A+\varphi_x)\,dx=\det(A)|\Omega|.
\tag{29}
\end{equation}

\Needspace{3\baselineskip}\textbf{基础知识}　行列式对矩阵元素的偏导为对应余子式；混合偏导交换；散度定理。约定
\((\varphi_x)_{ij}=\partial_j\varphi_i\)，\(\operatorname{cof}M\)
是余子式矩阵，满足
\((\operatorname{cof}M)_{ij}=\partial(\det M)/\partial M_{ij}\)。

\Needspace{3\baselineskip}\textbf{思路来源}　不要要求 \(A\) 可逆，也不要假设 \(Ax+\varphi(x)\)
是一一映射。引入连续扰动 \(A+t\varphi_x\)，证明其行列式积分对 \(t\)
的导数为零。关键是把导数写成散度，余子式自身的散度由混合偏导逐项抵消。

\Needspace{3\baselineskip}\textbf{证明}　先证明所需的余子式恒等式。设
\(u\in C^2(\Omega;\mathbb R^n)\)，记
\(C_{ij}=(\operatorname{cof}Du)_{ij}\)，则

\begin{equation}
\sum_{j=1}^n\partial_j C_{ij}=0,
\qquad i=1,\ldots,n.
\tag{30}
\end{equation}

当 \(n=1\) 时余子式为 \(1\)，结论显然。对 \(n\ge2\)，按排列展开，

\begin{equation}
C_{ij}=\sum_{\substack{\sigma\in S_n\\\sigma(i)=j}}
\operatorname{sgn}(\sigma)\prod_{k\ne i}\partial_{\sigma(k)}u_k.
\tag{31}
\end{equation}

对（31）求 \(\partial_j\)，并对 \(j\) 求和，每项由一个排列 \(\sigma\)
及一个 \(l\ne i\) 标记，其值为

\begin{equation}
\operatorname{sgn}(\sigma)
\big(\partial_{\sigma(i)}\partial_{\sigma(l)}u_l\big)
\prod_{k\ne i,l}\partial_{\sigma(k)}u_k.
\tag{32}
\end{equation}

将 \(\sigma(i)\) 与 \(\sigma(l)\) 对换，得到排列 \(\tau\)。对应于
\((\tau,l)\)
的乘积与（32）相同，因为混合偏导交换，其余因子不变；而排列符号相反。这一配对没有不动项，故所有项两两抵消，（30）得证。它通常称为
Piola 恒等式；这里已给出其直接证明。

令 \(u_t(x)=Ax+t\varphi(x)\)，\(M_t=Du_t=A+tD\varphi\)，并设

\begin{equation}
J(t)=\int_\Omega\det M_t\,dx.
\tag{33}
\end{equation}

行列式是 \(t\) 的次数不超过 \(n\)
的多项式，各系数可积，所以可直接逐项微分。行列式的微分公式和（30）给出

\begin{equation}
\begin{aligned}
J'(t)&=\int_\Omega\sum_{i,j}C_{ij}(t)\partial_j\varphi_i\,dx,\\
\sum_{i,j}C_{ij}(t)\partial_j\varphi_i
&=\sum_j\partial_j\left(\sum_i\varphi_i C_{ij}(t)\right)
-\sum_i\varphi_i\sum_j\partial_j C_{ij}(t)\\
&=\operatorname{div}V_t,
\quad (V_t)_j=\sum_i\varphi_i C_{ij}(t).
\end{aligned}
\tag{34}
\end{equation}

在 \(C_0^2\) 的紧支解释下，\(V_t\) 也紧支于
\(\Omega\)，其散度积分为零；这可把 \(V_t\)
延零至全空间后逐坐标积分验证。若采用
\(\varphi\in C^2(\overline\Omega)\)、边界取零的解释，则 \(V_t=0\)
于边界，Gauss 公式同样给出

\begin{equation}
J'(t)=\int_{\partial\Omega}V_t\cdot\nu\,dS=0.
\tag{35}
\end{equation}

因此 \(J\)
为常数，\(J(1)=J(0)=\det(A)|\Omega|\)，即（29）。证明只使用多项式微分，因而允许
\(A\) 或某些 \(M_t\) 奇异，也不需要 \(\det M_t\) 保持符号。

\Needspace{3\baselineskip}\textbf{易错条件}　余子式矩阵与伴随矩阵互为转置，式（30）的求和方向必须与
\(Du\) 的下标约定相符。若把行列式换成绝对值，结论一般不成立：例如
\(n=1\)、\(A=0\)、\(\Omega=(0,1)\)，任取非零
\(\varphi\in C_c^2(0,1)\)，则 \(\int_0^1\varphi'=0\)，但
\(\int_0^1|\varphi'|>0\)。

\Needspace{3\baselineskip}\textbf{方法总结}　将待证恒等式嵌入一个参数族，证明沿参数的导数为边界项，再用零边界消去它。行列式的反对称性使余子式散度自动消失，是本题中真正起作用的代数结构。

\Needspace{9\baselineskip}
\section{问题 B295}\label{ux95eeux9898-b295}

\textbf{原题号：} 习题 11.5.B 第 3
题。\href{https://zj139.github.io/math-notes/files/problems/楼分析习题-v2.pdf\#page=38}{原题，PDF
第 38 页（印刷第 36 页）}

\textbf{题目}　设 \(n\ge1\)，\(R\) 是仅在 \((0,0)\) 为零的 \(2n\)
次二元多项式，\(P,Q\) 是不超过 \(2n-2\) 次的二元多项式。若

\begin{equation}
\lim_{(x,y)\to\infty}\frac{R(x,y)}{(x^2+y^2)^n}>0,
\qquad
\frac{\partial}{\partial x}\frac{Q(x,y)}{R(x,y)}
=\frac{\partial}{\partial y}\frac{P(x,y)}{R(x,y)}
\quad((x,y)\ne(0,0)),
\tag{36}
\end{equation}

证明曲线积分 \(\int_C(P\,dx+Q\,dy)/R\) 绕原点的循环常数为零。

\Needspace{3\baselineskip}\textbf{基础知识}　穿孔平面上的闭微分形式；环域上的 Green
公式；大圆长度与被积函数的一致估计。

\Needspace{3\baselineskip}\textbf{思路来源}　闭性保证不同半径圆周上的积分相等，但穿孔平面不是单连通区域，不能仅据此断言积分为零。把圆周向无穷远推移，利用次数差获得
\(O(r^{-2})\) 的系数衰减，再乘以 \(O(r)\) 的圆周长度，才能消去循环常数。

\Needspace{3\baselineskip}\textbf{证明}　令 \(U=P/R\)、\(V=Q/R\)。由 \(R\) 在原点外不为零，\(U,V\)
在穿孔平面上光滑。记 \(C_\rho\) 为逆时针圆周 \(x^2+y^2=\rho^2\)，令

\begin{equation}
c(\rho)=\int_{C_\rho}U\,dx+V\,dy.
\tag{37}
\end{equation}

对任意 \(0<\rho<r\)，在圆环 \(\rho^2<x^2+y^2<r^2\) 上应用 Green
公式，得到

\begin{equation}
c(r)-c(\rho)
=\iint_{\rho^2<x^2+y^2<r^2}(V_x-U_y)\,dx\,dy=0.
\tag{38}
\end{equation}

所以 \(c(\rho)\) 与半径无关，它就是所求循环常数。

由（36）的正极限，存在 \(r_0\ge1\) 及 \(c_0>0\)，使在
\(x^2+y^2\ge r_0^2\) 上，\(R(x,y)\ge c_0(x^2+y^2)^n\)。另一方面，由
\(P,Q\) 的次数不超过 \(2n-2\)，存在 \(K>0\)，使在任意 \(r\ge r_0\)
的圆周上

\begin{equation}
\begin{aligned}
|P(x,y)|+|Q(x,y)|&\le K r^{2n-2},\\
\sqrt{U(x,y)^2+V(x,y)^2}
&\le\frac{|P|+|Q|}{R}\le\frac{K}{c_0r^2}.
\end{aligned}
\tag{39}
\end{equation}

按切向量的 Cauchy--Schwarz 不等式，

\begin{equation}
|c(r)|\le\int_{C_r}\sqrt{U^2+V^2}\,ds
\le 2\pi r\frac{K}{c_0r^2}
=\frac{2\pi K}{c_0r}\longrightarrow0.
\tag{40}
\end{equation}

结合（38），对任意固定 \(\rho>0\) 都有 \(c(\rho)=0\)。若 \(C\)
为包围原点的逆时针分段 \(C^1\)
简单闭曲线，取包围它的大圆，在二者之间的区域应用同一 Green 公式，也得
\(\int_C U\,dx+V\,dy=0\)；这与循环常数的定义一致。

\Needspace{3\baselineskip}\textbf{易错条件}　二维无穷远极限提供的是对所有方向统一的下界；只沿某条射线得到增长阶不足以证明（39）。也不能在含原点的圆盘内直接应用
Green 公式。作为对照，\((-y\,dx+x\,dy)/(x^2+y^2)\)
在穿孔平面上闭合，但其圆周积分为 \(2\pi\)；它的分子次数为
\(1\)，不满足本题 \(n=1\) 时次数至多 \(0\) 的条件。

\Needspace{3\baselineskip}\textbf{方法总结}　闭性使周期在环域中保持不变，无穷远估计再迫使该不变量为零。证明需要同时使用拓扑上的路径变形与分析上的统一衰减，二者缺一不可。

\Needspace{9\baselineskip}
\section{七题方法总结}\label{ux4e03ux9898ux65b9ux6cd5ux603bux7ed3}

\begin{longtable}[]{@{}
  >{\raggedright\arraybackslash}p{(\columnwidth - 4\tabcolsep) * \real{0.12}}
  >{\raggedright\arraybackslash}p{(\columnwidth - 4\tabcolsep) * \real{0.44}}
  >{\raggedright\arraybackslash}p{(\columnwidth - 4\tabcolsep) * \real{0.44}}@{}}
\toprule\noalign{}
\begin{minipage}[b]{\linewidth}\raggedright
题号
\end{minipage} & \begin{minipage}[b]{\linewidth}\raggedright
核心结果
\end{minipage} & \begin{minipage}[b]{\linewidth}\raggedright
关键方法
\end{minipage} \\
\midrule\noalign{}
\endhead
\bottomrule\noalign{}
\endlastfoot
B289 & 定积分性质移植到有限长曲线，反向不变 &
弧长参数、正权积分、连通性与重数辨析 \\
B290 & 球面积分化为 \(2\pi\int_{-1}^1f(Au)\,du\) &
正交旋转、等高参数、先证绝对可积 \\
B291 & \(\lvert\int_C f\,dz\rvert\le\int_C\lvert f\rvert\,ds\)
及等号条件 & 固定相位、单位切向量、非负差积分 \\
B292 & \(\iint P\,dy\,dz=\iint P\nu_1\,dS\) & 全局定向、参数
Jacobian、可积性条件 \\
B293 & 椭球面积分恒为 \(4\pi\) & 法向量识别、反平方场、挖孔 Gauss
公式 \\
B294 & 行列式积分只由固定边界数据决定 &
参数扰动、余子式配对抵消、散度边界项 \\
B295 & 绕原点循环常数为零 & 环域 Green 公式、次数差、大圆统一估计 \\
\end{longtable}

\end{document}
